\section{snnTorch}
\texttt{snnTorch} is a Python-based framework designed to facilitate the development, training, and simulation of SNNs within the PyTorch ecosystem. It provides modular building blocks such as spiking neuron models, surrogate gradient methods, and event-driven learning rules, allowing researchers to integrate biologically inspired temporal dynamics directly into deep learning architectures. By leveraging PyTorch's automatic differentiation and GPU acceleration, \texttt{snnTorch} enables efficient experimentation with hybrid SNNs and ANNs models and supports both supervised and unsupervised learning paradigms. The library has become a valuable tool for studying energy-efficient computation and neuromorphic learning systems, bridging the gap between neuroscience-inspired models and modern machine learning frameworks \cite{snntorch_docs}.

\subsection{Core Framework Components}
The framework consists of four primary modules. The \textbf{spiking neuron library} provides neuron models including Leaky Integrate-and-Fire (LIF), Synaptic, and Alpha neurons, as well as recurrent variants (SLSTM, SConv2dLSTM) for temporal processing. The \textbf{surrogate gradient module} implements multiple gradient approximation methods, including fast sigmoid, arctangent, and custom user-defined functions, to overcome the non-differentiability of spike generation. The \textbf{spike generation and data conversion library} enables transformation of static datasets into spike trains and supports event-driven data formats. The \textbf{functional module} provides utilities for common operations including loss computation (rate-based loss, cross-entropy loss), regularization, spike visualization, and layer-wise normalization (batch normalization through time).

\subsection{Practical Architecture Patterns} 
snnTorch enables construction of hybrid architectures that seamlessly combine traditional convolutional and recurrent layers with spiking neurons. For example, a typical convolutional SNN might stack layers such as \texttt{nn.Conv2d} followed by \texttt{snn.Leaky} neurons with configured decay rates and surrogate gradients, culminating in output-layer membrane potential extraction for continuous-valued predictions. Unsupervised learning is supported through spike autoencoders, where encoder-decoder structures integrate spiking neurons at multiple levels, and the framework's state management utilities (\texttt{utils.reset}) handle neuron membrane potential resets across time steps. The modular design allows researchers to experiment with different surrogate gradient functions and hyperparameter sweeps (decay rate 
β, membrane threshold) to optimize for both accuracy and energy efficiency.



